\newpage
\phantomsection
\label{prf:archimedean-q}
\LRAProofFor{thm:archimedean-q}

\begin{remark*}[Return]
\hyperref[thm:archimedean-q]{Return to Theorem}
\end{remark*}

\begin{theorem*}[Archimedean Property of $\mathbb{Q}$]
For every $x\in\mathbb{Q}$ there exists $n\in\mathbb{N}$ such that
\[
n>x.
\]
Equivalently, for every $x>0$ there exists $n\in\mathbb{N}$ such that
\[
\frac{1}{n}<x.
\]
\end{theorem*}

\begin{proof}
\textbf{Professional Standard Proof.}~\\
Let $x \in \mathbb{Q}$. Then there exist $a,b \in \mathbb{Z}$ with
$b \neq 0$ such that
\[
x=\frac{a}{b}.
\]
Replacing $a$ and $b$ by $-a$ and $-b$ if necessary, assume that
$b>0$.

If $a<0$, then
\[
x=\frac{a}{b}<0<1,
\]
so $n=1$ satisfies $n>x$.

Assume instead that $a\ge 0$. Let
\[
n=a+1.
\]
Then $n\in\mathbb{N}$. Since $b>0$ and $b\in\mathbb{Z}$, one has
$b\ge 1$. Hence
\[
b(a+1)=ab+b\ge a+b>a.
\]
Dividing by the positive integer $b$ gives
\[
a+1>\frac{a}{b}.
\]
Therefore
\[
n>x.
\]
Thus, for every $x\in\mathbb{Q}$, there exists $n\in\mathbb{N}$ such
that $n>x$.

Now let $x\in\mathbb{Q}$ with $x>0$. Since $1/x\in\mathbb{Q}$, the
first part gives some $n\in\mathbb{N}$ such that
\[
n>\frac{1}{x}.
\]
Because $x>0$ and $n>0$, multiplication by $x$ and division by $n$
preserve the order, so
\[
x>\frac{1}{n}.
\]
Hence
\[
\frac{1}{n}<x.
\]
\end{proof}

\begin{proof}
\textbf{Detailed Learning Proof.}~\\

\textbf{Step 1.} Represent the rational number with positive denominator.

Fix $x\in\mathbb{Q}$. By the definition of rational number, there exist
integers $a,b\in\mathbb{Z}$ with $b\neq 0$ such that
\[
x=\frac{a}{b}.
\]
If $b<0$, replace $a$ and $b$ by $-a$ and $-b$. This does not change
the rational number represented, and it gives a representation
\[
x=\frac{a}{b}
\]
with
\[
b>0.
\]

\textbf{Step 2.} Handle the case where the rational number is negative.

If $a<0$, then since $b>0$,
\[
\frac{a}{b}<0.
\]
Therefore
\[
x<0<1.
\]
Since $1\in\mathbb{N}$, the choice
\[
n=1
\]
satisfies
\[
n>x.
\]

\textbf{Step 3.} Handle the case where the rational number is nonnegative.

Assume $a\ge 0$. Define
\[
n=a+1.
\]
Then $n\in\mathbb{N}$. Since $b\in\mathbb{Z}$ and $b>0$, it follows
that
\[
b\ge 1.
\]
Consequently,
\[
b(a+1)=ab+b\ge a+b>a.
\]
Since $b>0$, division by $b$ preserves the strict inequality:
\[
a+1>\frac{a}{b}.
\]
Thus
\[
n>x.
\]
This proves that for every rational number $x$ there exists
$n\in\mathbb{N}$ such that
\[
n>x.
\]

\textbf{Step 4.} Derive the reciprocal formulation.

Let $x\in\mathbb{Q}$ and suppose $x>0$. Then
\[
\frac{1}{x}\in\mathbb{Q}.
\]
Applying the first part to $1/x$, there exists $n\in\mathbb{N}$ such
that
\[
n>\frac{1}{x}.
\]
Multiplying by the positive rational number $x$ gives
\[
nx>1.
\]
Dividing by the positive natural number $n$ gives
\[
x>\frac{1}{n}.
\]
Therefore
\[
\frac{1}{n}<x.
\]
\end{proof}

\begin{remark*}[Proof structure]
The argument is direct. The first formulation is proved by writing an
arbitrary rational number as $a/b$ with positive denominator and then
choosing an integer strictly larger than that quotient. The reciprocal
formulation is then obtained by applying the first formulation to $1/x$
when $x>0$.
\end{remark*}

\begin{dependencies}
\begin{itemize}
  \item \hyperref[def:rational-numbers]{Definition~\ref*{def:rational-numbers}}
  \item Basic arithmetic of integers and rationals
  \item Order preservation under multiplication and division by positive rational numbers
\end{itemize}
\end{dependencies}
